% =====================================================================
%  Results and Discussion — conference paper section (IEEE two-column)
%  Source of numbers: docs/methodology_and_results.md (as recorded) and
%  results/exp19_sequence_coverage.json (committed).
%
%  Requires in the preamble:
%    \usepackage{booktabs,multirow,graphicx,xcolor}
%    \newcommand{\stale}{\textsuperscript{\textcolor{red}{\dag}}}
%  Every number marked \stale comes from a run that predates a code change
%  (exp9, exp14, exp15, exp17, exp18) and MUST be re-run before submission;
%  delete the macro and the footnote in Sec. V-A once refreshed.
%  Figure paths assume figures are copied next to the .tex (docs/figures/).
% =====================================================================

\section{Results}
\label{sec:results}

\subsection{Experimental setup}
\label{sec:setup}

All experiments use a single emitter on $M=16$ channels and a slot-synchronous
detection-level sensor: the true channel is detected with probability $P_d$ and
every other channel raises a false alarm independently with probability $P_{fa}$.
The generator library contains the in-bank families (LFSR, filtered LFSR, LCG,
full-state LCG, weak-seed PRNG, first-order Markov, logistic/tent/H\'enon maps),
three \emph{held-out} families for which no expert exists (lagged-Fibonacci,
quadratic congruential, sine map), eight structured sequence types, and a keyed
SHA-256 CSPRNG as the unpredictable control. Experts see only observations; true
channels and generator parameters are never passed to the bank (information
parity, Sec.~\ref{sec:method}).

Unless stated otherwise, identification uses MSPRT with $\delta=0.05$ and the
conference bank (one hypothesis per family plus $H_0$), windows of $T=100$--$160$
slots, and 30--50 seeds per condition. We report identification accuracy,
hops-to-identify $\tau$ over correct declarations, and two error rates: any
false identification, and \emph{hard} false identification, i.e.\ declaring an
algebraic family (LFSR, filtered LFSR, LCG, weak-seed) that is not the truth.
Prediction is scored by top-1 next-channel accuracy against the chance line
$1/M=0.0625$; ``steady state'' is the final 40 slots. Intervals are bootstrap
95\% CIs; paired comparisons use a paired bootstrap and McNemar's test.
Values marked \stale{} predate the final code revision and are being re-run.

% ---------------------------------------------------------------------
\subsection{Identification speed and the stopping law}
\label{sec:res-speed}

Table~\ref{tab:hops} reports hops-to-identify as the sensor thins. The three
deterministic in-bank families are identified at 100\% accuracy in 4--5 hops
with a perfect sensor and in 11--13 hops at $P_d=0.5$. The Markov hopper needs
roughly ten times as many hops (46--74) and its accuracy falls to 0.83 at
$P_d=0.5$, because its transition matrix has to be estimated online before it
predicts better than the uniform $H_0$.

\begin{table}[t]
\centering
\caption{Mean hops-to-identify versus $P_d$ (MSPRT, $\delta=0.05$,
$P_{fa}=0.02$, 30 seeds).}
\label{tab:hops}
\begin{tabular}{lcccccc}
\toprule
$P_d$ & 1.0 & 0.9 & 0.8 & 0.7 & 0.6 & 0.5 \\
\midrule
LFSR      & 4.2  & 4.7  & 5.5  & 7.1  & 8.3  & 10.9 \\
Weak-seed & 4.2  & 4.8  & 5.6  & 7.2  & 8.8  & 11.6 \\
LCG       & 4.9  & 5.5  & 6.7  & 8.1  & 10.0 & 13.3 \\
Markov    & 45.9 & 52.5 & 55.5 & 65.9 & 68.1 & 74.3 \\
\bottomrule
\end{tabular}
\end{table}

The growth follows the sequential-test law
$\mathbb{E}[\tau]\approx\log(1/\delta)/(P_d\,\min_{k\neq k^*}\mathrm{KL})$.
On a Markov source, chosen because its per-slot KL is finite and the law is
therefore visible, $\mathbb{E}[\tau]$ is linear in $\log(1/\delta)$
($R^2=0.984$) and in $1/P_d$ ($R^2=0.997$). This confirms that a missed
detection removes one likelihood factor and otherwise does no damage. The
deterministic families saturate the law: one inconsistent observation removes a
wrong candidate, so $\delta$ hardly matters for them.

% ---------------------------------------------------------------------
\subsection{The value of a sensor-aware likelihood}
\label{sec:res-sensor}

Table~\ref{tab:ablation} isolates the emission model. The same bank is run with
the sensor-aware likelihood and with a naive clean-match score at a stressing
$P_{fa}=0.05$. The two agree for $P_d\ge0.8$. Below that the naive bank slows
down and starts naming wrong generators: at $P_d=0.5$ its accuracy is 0.60
against 0.98, and it needs 21.1 hops against 15.3. The naive score counts a
missed hop as evidence against the true generator and a phantom as evidence for
a wrong one. The sensor-aware emission does neither, since a phantom raises
every candidate by the same factor and cancels in the posterior.

\begin{table}[t]
\centering
\caption{Sensor-aware vs naive likelihood ($P_{fa}=0.05$, 30 seeds per generator).}
\label{tab:ablation}
\begin{tabular}{llcccccc}
\toprule
& $P_d$ & 1.0 & 0.9 & 0.8 & 0.7 & 0.6 & 0.5 \\
\midrule
\multirow{2}{*}{Accuracy} & aware & 1.00 & 1.00 & 1.00 & 0.99 & 0.98 & \textbf{0.98} \\
                          & naive & 1.00 & 1.00 & 1.00 & 0.99 & 0.90 & 0.60 \\
\midrule
\multirow{2}{*}{Hops}     & aware & 4.7 & 5.6 & 7.2 & 8.9  & 11.1 & \textbf{15.3} \\
                          & naive & 4.7 & 5.5 & 8.0 & 13.4 & 16.7 & 21.1 \\
\bottomrule
\end{tabular}
\end{table}

The same effect is sharper at the level of a single recovery algorithm. For an
$L=7$ LFSR, plain Berlekamp--Massey recovers the register on every run at
$P_d=1$ and on \emph{no} run for any $P_d<1$, because it needs about $2L$
consecutive clean outputs and one miss breaks the run. State-enumeration
recovery under the detection channel keeps a posterior over register states
through gaps and recovers on every run across the whole $P_d$ sweep. False
alarms are equally harmless to the in-bank families: with $P_d=0.8$ and
$P_{fa}$ up to 0.1, accuracy stays at 1.00 and false identification at 0.00.

% ---------------------------------------------------------------------
\subsection{Open-set safety}
\label{sec:res-safety}

\emph{None-of-the-above.} On the keyed CSPRNG the bank declares
\emph{unknown} in 87\% of runs and never claims an algebraic break
(hard false-ID 0/40, 95\% upper bound 7.5\%). On a filtered LFSR, which the
conference bank has no expert for, it falls back to the Markov hypothesis, a
soft error, again with no hard false identification. A pairwise SPRT without
$H_0$ names a generator on every such run (100\% false-ID). Any closed-set test
must do this, so the comparison shows the need for an explicit unknown class
rather than a weakness peculiar to SPRT.

\emph{Stale reconnaissance.} When the secret is placed outside the searched
range (LCG seed beyond the searched keyspace, weak-seed keyspace $2^{16}$
against 256 searched, LFSR length 10 against $\le7$), hard false-ID is
0.00/0.03/0.03 and the bank routes to unknown in 63\%/87\%/30\% of runs. With no
in-range candidate, every deterministic likelihood stays near uniform, so the
bank falls back to $H_0$ or a soft Markov answer instead of asserting a break.

\emph{Method comparison.} Table~\ref{tab:methods} compares decision rules on
identical episodes and likelihoods. The bank is the most accurate method and
the only one that is both accurate and low in false identification. The
fixed-length and SPRT tests lose accuracy on exactly the out-of-bank episodes
where they are forced to name a class. The learned classifier has an
\emph{unknown} output and therefore a low false-ID rate, but summary features
cannot separate an LCG from a weak-seed PRNG from a CSPRNG, so its accuracy is
0.38.

\begin{table}[t]
\centering
\caption{Decision rules on the same mixed in/out-of-bank episodes
($P_d=0.8$, $P_{fa}=0.02$, 30 seeds). $\Delta$: paired-bootstrap accuracy gain
of the bank; all $p<10^{-3}$.}
\label{tab:methods}
\begin{tabular}{lcccc}
\toprule
Method & Acc. & False-ID & Hops & $\Delta$ \\
\midrule
\textbf{Bank + MSPRT (ours)} & \textbf{0.79} & 0.19 & 21.1 & --- \\
Pairwise SPRT (no $H_0$)     & 0.66 & 0.34 & 5.8  & +0.139 \\
Fixed-length LRT             & 0.67 & 0.33 & 30.0 & +0.128 \\
Learned classifier           & 0.38 & 0.11 & 2.0  & +0.411 \\
\bottomrule
\end{tabular}
\end{table}

\emph{Coverage at scale.} Table~\ref{tab:coverage} aggregates 7{,}500 episodes
(10 families $\times$ 5 $P_d$ $\times$ 3 $P_{fa}$ $\times$ 50 seeds) through the
one-expert-per-family identification bank. Accuracy is at least 0.92 for nine of
the ten families. Hard false-ID is at most 0.01 for nine families and 0.05 for
the tent map, whose quantised dynamics sometimes admit a short consistent
deterministic fit. The full-state LCG is the exception in accuracy (0.50). At
$M=16$ its period is at most 16, so it is also an exact Markov cycle; when $P_d$
is too low to supply the six consecutive clean outputs the cryptanalytic expert
needs, the Markov expert legitimately wins. This is a soft error, and accuracy
is 0.975 at $P_d=0.95$. With the full 13-expert bank at $P_d=0.95$, the
confusion matrix is diagonal: each generator is won by its intended expert,
plain BM never beats its gap-tolerant counterpart, and the CSPRNG goes to $H_0$
(0.925).

\begin{table}[t]
\centering
\caption{Per-family identification over 7{,}500 episodes
($P_d\in[0.6,1]$, $P_{fa}\in\{0,0.02,0.05\}$, 50 seeds per cell).}
\label{tab:coverage}
\begin{tabular}{lcc|lcc}
\toprule
Family & Acc. & Hard FID & Family & Acc. & Hard FID \\
\midrule
LFSR          & 0.99 & 0.01 & Logistic & 1.00 & 0.00 \\
Filtered LFSR & 1.00 & 0.00 & Tent     & 0.95 & 0.05 \\
LCG           & 0.99 & 0.01 & H\'enon  & 0.99 & 0.01 \\
Full-state LCG& 0.50 & 0.01 & Markov   & 0.92 & 0.01 \\
Weak-seed     & 1.00 & 0.00 & CSPRNG   & 0.88 & 0.01 \\
\bottomrule
\end{tabular}
\end{table}

% ---------------------------------------------------------------------
\subsection{Anytime-valid stopping and calibration}
\label{sec:res-guarantee}

With e-process stopping, the probability of ever declaring a wrong in-bank
generator is bounded by Ville's inequality. Sweeping
$\beta\in\{0.5,\dots,0.01\}$ on in-bank truths ($P_d=0.7$, $P_{fa}=0.05$,
$\ge100$ seeds per family), the empirical false-ID is 0.24, 0.10, 0.03, 0.027,
0.01 and 0.01\stale{}, on or below $y=\beta$ at every point. The bound is
conservative here because the deterministic families separate quickly. We
neither claim nor test the bound for out-of-bank truths, where the martingale
argument does not apply and protection comes from $H_0$
(Sec.~\ref{sec:res-safety}).

The identification posterior is well calibrated in the high-confidence bin,
where most declarations fall and accuracy is close to 1. The expected
calibration error over 180 declarations is 0.154. Almost all of it comes from
boundary and out-of-bank sources (Markov, CSPRNG, filtered LFSR), where a short
window sometimes produces confident mid-range errors.

% ---------------------------------------------------------------------
\subsection{Next-channel prediction}
\label{sec:res-predict}

Table~\ref{tab:predict} reports the headline metric. Once a deterministic
generator is identified, its future is known, so the predictor reaches 1.00
top-1 on LFSR, weak-seed and LCG emitters. The bank's own next-hop distribution
already reaches log-loss 0.07--0.09, against $\log M=2.77$ for uniform. On the
CSPRNG both predictors fall to chance (0.07), which is the correct answer.

The held-out families test whether prediction outlives a limited bank. On
lagged-Fibonacci, whose next channel depends on the previous two, the Fixed-Share
predictor with the training-free universal expert reaches 0.88 top-1, against
0.42 for the identification-only mixture. The bank's best available model is
first-order and cannot represent the rule. Across a $P_d$ sweep from 1.0 to 0.7
the proposed predictor falls from 0.96 to 0.40. The bank-only predictor is
flat at 0.34--0.40, so the gap is large at high $P_d$ and closes by
$P_d=0.7$.\stale{} The quadratic congruential generator is
predicted well by both (0.95/0.99), because its full-period order-1 cycle is
also captured by the Markov expert. The chaotic sine map is predicted at about
seven times chance by both (0.43/0.44).

\begin{table}[t]
\centering
\caption{Steady-state top-1 next-channel accuracy ($P_d=0.95$, $P_{fa}=0$,
$T=120$, 12 seeds; chance $=0.0625$). Held-out families have no expert.}
\label{tab:predict}
\begin{tabular}{lcccc}
\toprule
Generator & Proposed & Bank only & Markov & Uniform \\
\midrule
LFSR        & 1.00 & 1.00 & 0.04 & 0.06 \\
Weak-seed   & 1.00 & 1.00 & 0.07 & 0.06 \\
LCG         & 1.00 & 1.00 & 0.07 & 0.06 \\
Logistic    & 0.54 & 0.54 & 0.06 & 0.06 \\
Markov      & 0.38 & 0.37 & 0.10 & 0.06 \\
\midrule
\emph{Lagged-Fib.} (held out)   & \textbf{0.88}\stale & 0.42\stale & 0.26 & 0.06 \\
\emph{Quadratic CG} (held out)  & 0.95\stale & 0.99\stale & 0.03 & 0.06 \\
\emph{Sine map} (held out)      & 0.43\stale & 0.44\stale & 0.05 & 0.06 \\
\midrule
CSPRNG (control) & 0.07 & 0.05 & 0.06 & 0.06 \\
\bottomrule
\end{tabular}
\end{table}

\emph{Structured sequences and secrets outside reconnaissance.}
Table~\ref{tab:structured} adds the structured experts (Bayesian rule search,
full-keyspace seed beam, adaptive-window BM) and scores every structured
sequence type on a clean and a lossy sensor. Every type moves to or near its
ceiling. Hops-to-lock (the median observed slots after which the next five
predictions are all correct) are close to the information minimum: 4 hops for an
arithmetic progression, $d+3$ for a degree-$d$ polynomial, and just over one
period for a lookup table, which cannot be predicted before one period has
been seen. Seed search over the whole keyspace locks a weak-seed PRNG with a
$2^{18}$ keyspace in 5--7 hops, where the range-limited expert never locks. Two
cases stay out of reach, as they should. A 16-bit LFSR is recovered on a clean
stream (1.00) but not on a lossy one (0.12), because Berlekamp--Massey needs
about 11 consecutive clean hops. The CSPRNG stays at chance.

\begin{table}[t]
\centering
\caption{Steady-state top-1 before $\to$ after the structured experts, and
median hops-to-lock (after), $T=120$, 8 seeds.}
\label{tab:structured}
\setlength{\tabcolsep}{3.5pt}
\begin{tabular}{lcc|cc}
\toprule
& \multicolumn{2}{c|}{$P_d{=}0.95,\,P_{fa}{=}0$} & \multicolumn{2}{c}{$P_d{=}0.85,\,P_{fa}{=}0.02$} \\
Sequence type & top-1 & lock & top-1 & lock \\
\midrule
Arithmetic            & 0.98$\to$\textbf{1.00} & 4    & 0.86$\to$\textbf{0.99} & 6 \\
Polynomial ($d{=}2$)  & 0.93$\to$\textbf{1.00} & 5    & 0.62$\to$\textbf{1.00} & 9 \\
Linear rec.\ (ord.\ 3)& 0.72$\to$\textbf{1.00} & 7.5  & 0.42$\to$\textbf{0.98} & 14 \\
Bilinear rec.         & 0.86$\to$\textbf{1.00} & 6.5  & 0.66$\to$\textbf{0.99} & 11 \\
Lookup table          & 0.86$\to$\textbf{1.00} & 22.5 & 0.40$\to$\textbf{0.95} & 36.5 \\
Table + drift         & 0.37$\to$\textbf{1.00} & 28.5 & 0.10$\to$\textbf{0.92} & 38 \\
Rule switch / 40 hops & 0.50$\to$\textbf{0.84} & 7    & 0.27$\to$\textbf{0.71} & 12.5 \\
Noisy rule (10\%)     & 0.84$\to$\textbf{0.87} & 14.5 & 0.62$\to$\textbf{0.72} & 13.5 \\
\midrule
Weak seed, $2^{18}$   & 0.07$\to$\textbf{1.00} & 5    & 0.07$\to$\textbf{1.00} & 7 \\
LCG, any seed         & 0.20$\to$\textbf{1.00} & 4    & 0.18$\to$\textbf{1.00} & 6 \\
LFSR, $L{=}16$        & 0.06$\to$\textbf{1.00} & 15.5 & 0.04$\to$0.12          & 61 \\
\midrule
CSPRNG (control)      & 0.06$\to$0.07          & ---  & 0.06$\to$0.07          & --- \\
\bottomrule
\end{tabular}
\end{table}

% ---------------------------------------------------------------------
\subsection{Tracking a generator switch}
\label{sec:res-switch}

Table~\ref{tab:switch} replays random switch plans over all 13 algorithms, with
regimes of 50--90 hops and a fresh secret at every switch, through the
switch tracker. Switches are detected in 90--100\% of cases. The median delay is
0 hops on a perfect sensor and 11 on a contested one. False alarms stay at or
below 0.25 per 100 hops. Every switch \emph{into} a family the bank has an
expert for was detected on every sensor. The misses are between two sources
with no exact structure (e.g.\ Markov $\to$ CSPRNG), where the old segment's
learners simply adapt. From hop 6 after a switch, next-hop top-1 is 71--79\% on
perfect and strong sensors, matching an oracle restarted at the true switch.
Hops 1--5 are the evidence any detector needs to notice the change. One lock
in 36 regimes on the contested sensor was dangerous: an LFSR candidate fitted a
stretch of a Markov hopper's noisy hops, mostly through ghost detections.

\begin{table}[t]
\centering
\caption{Switch tracking: 6 random plans $\times$ 400 hops per sensor, 13
algorithms, $M=16$.\stale}
\label{tab:switch}
\setlength{\tabcolsep}{3pt}
\begin{tabular}{lcccc}
\toprule
Sensor ($P_d$, $P_{fa}$) & Perfect & Strong & Typical & Contested \\
& (1, 0) & (.95, .01) & (.8, .02) & (.6, .05) \\
\midrule
Switches detected      & 100\% & 97\% & 93\% & 90\% \\
Delay, median/90th pct & 0/6 & 1/6 & 4/12 & 11/24 \\
False alarms /100 hops & 0.04 & 0.04 & 0.25 & 0.13 \\
Named exactly at end   & 81\% & 72\% & 78\% & 75\% \\
Dangerous locks        & 0\% & 0\% & 0\% & 3\% \\
\midrule
Top-1, hops 1--5 after & 31\% & 21\% & 9\%  & 5\% \\
Top-1, hops 6--10      & 79\% & 71\% & 55\% & 21\% \\
Top-1, hops 21--40     & 77\% & 75\% & 72\% & 56\% \\
\bottomrule
\end{tabular}
\end{table}

The response to an alarm also matters. On a near-deterministic Markov hopper
that switches at mid-window, resetting the mixture weights \emph{and} flushing
the universal predictor's context recovers top-1 prediction 1.2--2.3$\times$
faster than resetting the weights alone\stale{} (15--46 against 33--56 hops
over four concentration levels). The cumulative-evidence bank alone is slowest.

% ---------------------------------------------------------------------
% OPTIONAL (keep only if the paper retains the downstream-use claim; see
% docs/paper/experiment_review.md §3 item 14 before keeping it).
\subsection{Downstream use: acting on the predictive distribution}
\label{sec:res-alloc}

Feeding the whole predictive distribution, not just the top-1 channel, into a
threshold-aware greedy knapsack over channels gives spot and barrage jamming as
two ends of one rule. Against a CSPRNG the knapsack keeps a barrage-like floor
(reward 0.11) where water-filling, which spreads every channel below the
denial threshold, achieves 0.00. The two tie on a fully predictable LFSR
(0.94/0.96). Against the unpredictable CSPRNG, a reactive follower that jams
the currently observed channel raises reward from 0.11 at dwell 1 to 0.82 at
dwell 8: an unpredictable source can still be jammed if it dwells longer than
the reaction latency.

% =====================================================================
\section{Discussion}
\label{sec:discussion}

\emph{One likelihood explains most of the curves.} Every deterministic family
is scored the same way, by enumerating its parameters and weighting each
candidate trajectory through the reduced detection-level emission. This one
design choice accounts for three separate observations: single-digit hops to
identify (one inconsistent observation removes a candidate), $\tau$ that
scales with $1/P_d$ and $\log(1/\delta)$ (a miss removes a factor and leaves
the rest of the evidence intact), and immunity to false alarms (a phantom
scales every candidate equally). The ablation in Table~\ref{tab:ablation} and
the BM cliff measure the same thing at two levels. An algorithm that assumes a
clean stream fails at the first miss; one that models the sensor degrades only
as fast as the evidence does.

\emph{Hard versus soft errors.} The bank's \emph{any-wrong} error rate is
non-zero on boundary sources: a filtered LFSR read as Markov by the conference
bank, or a full-state LCG read as Markov at low $P_d$. These are soft errors.
The predicted next hop is still useful, and no false structural claim is made.
The dangerous error is a confident claim of an algebraic break that does not
exist, because it would lead an operator to commit resources to the wrong
generator. That rate is at most 0.01 for nine of ten families over 7{,}500
episodes, 0.05 for the tent map, and 1 of 36 regimes for the switch tracker
under a contested sensor. A closed-set test makes the same error on every
out-of-bank source. The case for an explicit generative $H_0$ is this
asymmetry between error types, not raw accuracy.

\emph{The predictability boundary is a result.} On a keyed CSPRNG the system
answers \emph{unknown}, its predictive distribution flattens, and top-1 falls
to $1/M$. A system that kept finding structure there would be overfitting
noise. Detecting the boundary quickly and saying so is the desired behaviour,
and it is what lets a downstream controller switch from prediction to reactive
or barrage tactics (Sec.~\ref{sec:res-alloc}).

\emph{A limited bank does not limit prediction as much.} Any bank of experts
is finite, and an adversary can always choose a generator outside it. The
Fixed-Share mixture weights experts by recent predictive skill, so when no
identification expert fits, the training-free universal expert takes the
weight within a few slots. The lagged-Fibonacci result (0.88 against 0.42) is
the clearest case. The sine-map and Markov rows show the limit: if a source's
structure is only statistical, the mixture can do no better than its best
statistical learner. Program search over structured rules narrows the gap
further (Table~\ref{tab:structured}), at a cost: a period-$P$ table cannot be
predicted before one period has been observed, and the measured lock times sit
close to that floor.

\emph{Switching adversaries.} An identifier that accumulates evidence is, by
construction, slow to forget. Two mechanisms address this at different
granularities. Fixed-Share with a context flush handles drift and isolated
switches cheaply. The switch tracker, whose run-length hypotheses each carry a
bank anchored at the hypothesised switch slot, handles re-keyed counter-based
generators, which an unanchored bank cannot re-identify at all. The
contested-sensor row of Table~\ref{tab:switch} shows the price: a detection
delay of about 11 hops and one dangerous lock, which matches the bank's own
residual hard false-ID rate under that sensor.

\emph{Limitations and threats to validity.} (i) All results are simulated.
The sensor is slot-synchronous with i.i.d.\ misses and is \emph{matched}: the
same $P_d$, $P_{fa}$ generate and score the data. Robustness to a mis-specified
or bursty sensor, hop-timing jitter and interleaved emitters is not yet
evaluated, and every parameter range is an assumption rather than a value
calibrated on recordings. (ii) $M=16$ throughout. The full-state LCG result
is confounded by its short period at this $M$, and runtime grows with the
expert set: the switch tracker costs 40--90\,ms per hop, which limits it to
slow hop rates without further engineering. (iii) The generators and the
experts were written by the same authors. The held-out families reduce, but do
not remove, the risk that the library favours the bank. (iv) The e-process
guarantee holds only for in-bank truths, and protection out of bank is
empirical. (v) Several headline prediction and tracking results use small seed
counts (6--12), and zero error rates over 40 seeds have a 95\% upper bound
near 7.5\%. These are the priorities for the journal version, together with
an external learned predictor baseline and a calibrated next-hop
distribution.
